\section[Finite-Enclosure Method]{Finite-Enclosure Obstructions}
\label{sec:finite-enclosure}

The finite-enclosure method identifies a compact set that the C triangle
must contain, then proves that this set does not fit inside any open unit
equilateral triangle.

Put
\[
 h:=\frac{\sqrt3}{2}.
\]

The closed-disk notation and $\Lambda$ are defined in
Section~\ref{sec:introduction}. By \zcref{lem:compact-cover-margin}, a
compact set contained in an open unit triangle has enclosure side below
one. Conversely, $\Lambda(K)<1$ gives such an open unit enclosure by
slightly enlarging a smaller closed enclosure. Thus it is enough to
exclude every open unit candidate containing the same fixed set $K$.

\subsection{The two tasks in a finite-enclosure proof}
Fix the original seven triangles and call their arrangement $\mathcal U$.
A nonempty compact witness set $K=K(\mathcal U)$ is then chosen once and for all. There are
two separate tasks:
\[
 \underbrace{K\subset U_C}_{\text{forced containment}}
 \qquad\hbox{and}\qquad
 \underbrace{\Lambda(K)\ge1}_{\text{enclosure obstruction}}.
\]
The first makes $\Lambda(K)<1$ by the compact covering margin, contradicting
the second. Most witnesses are points missed by all open V triangles.
A structural anchor, such as the C midpoint, may instead be included because
it is already known to belong to $U_C$; such an anchor need not be missed
by every V triangle.

To prove the second task, suppose an \emph{arbitrary} open unit triangle
$U'$ contains the fixed $K$. Its orientation, offsets, and radial exits may
vary. The original reaches and the witness coordinates do not. In particular,
$U'$ need not cover the original boundary gaps or complete the original
skeleton cover. Any inequality obtained from original coverage must be
established before $U'$ is introduced and retained as a fixed input.

\begin{center}
\small
\renewcommand{\arraystretch}{1.15}
\begin{tabular}{@{}p{.27\linewidth}p{.65\linewidth}@{}}
\toprule
Object & What is fixed or required \\
\midrule
Original arrangement $\mathcal U$ & Its actual reaches and all coverage
hypotheses of the current case. \\
Witness set $K(\mathcal U)$ & Chosen from the original arrangement; never
recomputed from the candidate. \\
Candidate $U'$ & Required only to contain $K$; its own geometric parameters
are recomputed. \\
Replacement arrangement & A different operation, used later to preserve
skeleton coverage, not an enclosure candidate. \\
\bottomrule
\end{tabular}
\end{center}

The six-point construction uses actual gap endpoints and total radial
endpoints. The four-point construction uses a supported radial endpoint
and boundary anchors. The zero-gap construction instead uses uniform
capacity bounds to choose six radial points and three frontier points.
Their common feature is fixed containment, not a common coordinate formula.


\subsection{The enclosure calculus: finitely many orientations}
\label{sec:universal-calculus}

The next four results are the general enclosure tools, rather than
case-specific estimates. The gauge removes translation and scale; the
support-cell argument reduces the remaining orientation variable to a
finite list. They do not construct the witnesses or prove that those
witnesses are forced into the C triangle.



Let $\mathsf R$ denote counterclockwise rotation through $2\pi/3$.  For a
nonempty compact $K\subset\mathbb R^2$, put
\[
 h_K(n):=\max_{x\in K}\langle x,n\rangle.
\]

\begin{proposition}[Equilateral enclosure gauge]
\label{prop:new-enclosure-gauge}
\label{prop:universal-enclosure-gauge}
The least side length $\Lambda(K)$ of a closed equilateral triangle containing
$K$ is
\begin{equation}
 \Lambda(K)=\frac2{\sqrt3}
 \min_{\lVert n\rVert=1}\sum_{j=0}^2h_K(\mathsf R^jn).
 \label{eq:universal-lambda}
\end{equation}
It is monotone under inclusion, translation invariant, and positively
homogeneous.
\end{proposition}

\begin{proof}
For fixed outward normals $n,\mathsf Rn,\mathsf R^2n$, the least containing
triangle has the corresponding three support numbers, and its side is
$2/\sqrt3$ times their sum.  Translation invariance follows from
\[
 n+\mathsf Rn+\mathsf R^2n=0,
\]
while monotonicity and homogeneity are immediate.
\end{proof}

\begin{lemma}[Support-cell rotation]
\label{lem:support-cell-rotation}
A least equilateral enclosure of a nondegenerate polygon can be chosen
with one side containing a polygon edge. For a centered disk together
with finitely many points, it suffices to test point--point and point--disk
ties, or a disk-only orientation. Here a \emph{support cell} is an open
angular interval on which the maximizing point or disk is fixed in each
of the three support directions; a \emph{tie} is an endpoint where a
maximizer changes.
\end{lemma}
\begin{proof}
By \zcref{prop:new-enclosure-gauge}, the side is the three-term support
sum divided by $h=\sqrt3/2$. Parametrize unit normals by
$n(\theta)=(\cos\theta,\sin\theta)$. For a disk of radius $\rho>0$, its
support value is $\rho$ in every unit direction. On a polygon cell with fixed support vertices the sum
is $g(\theta)=a\cos\theta+b\sin\theta>0$, so $g''=-g<0$ and no
interior minimum occurs. In the disk case a cell with $q$ disk supports
has sum $q\rho+\langle v,n(\theta)\rangle$, where $v$ is the sum
of the active point vectors rotated back by their respective multiples
of $120^\circ$. If a point is active its
projection exceeds $\rho\ge0$, and the inner product is positive, again
giving a negative second derivative. A disk-only cell has constant support sum $3\rho$. Each application must still identify its exposed ties.
\end{proof}

\subsubsection{Finite caliper criterion}
\label{subsec:finite-caliper}

Let $\mathsf J$ denote counterclockwise rotation through $\pi/2$.  For a
nonempty finite set $F$ define
\begin{equation}
 W_F(N):=\sum_{k=0}^2\max_{z\in F}
 \langle z,\mathsf R^kN\rangle.
 \label{eq:caliper-support-sum}
\end{equation}

\begin{theorem}[Finite caliper certificate]
\label{thm:cert-caliper}
If $F$ has at least two distinct points, then $F$ is contained in a closed
unit equilateral triangle if and only if there are distinct $p,q\in F$ and
$\varepsilon\in\{-1,1\}$ such that
\begin{equation}
 W_F\bigl(\varepsilon\mathsf J(q-p)\bigr)
 \le\frac{\sqrt3}{2}\lVert p-q\rVert.
 \label{eq:finite-caliper-test}
\end{equation}
A singleton is contained in equilateral triangles of arbitrarily small
positive side length.
\end{theorem}

\begin{proof}
Apply \zcref{prop:universal-enclosure-gauge,lem:support-cell-rotation}
to $\conv F$. A minimizing normal is perpendicular to a hull edge
$[p,q]$ and, after cyclic
rotation, equals
$\varepsilon\mathsf J(q-p)/\lVert p-q\rVert$.  Homogeneity gives
\eqref{eq:finite-caliper-test}; a nondegenerate segment follows by checking
its two normals.

\end{proof}

\begin{proposition}[Disk--finite-set calipers]
\label{prop:new-disk-finite-caliper}
Let $\rho>0$, $P=\{p_1,\ldots,p_m\}\subset\mathbb R^2$, and
$K=\conv(\mathbb B(O,\rho)\cup P)$. Use the rotation $\mathsf J$
defined above. Unless the disk is the active support in all three directions at a minimizing
orientation, a minimizing normal belongs to the finite list
\[
 \pm\frac{\mathsf J(p_k-p_i)}{\|p_k-p_i\|}
 \quad(p_i\ne p_k)
\]
or
\[
 \frac{\rho p_i\pm
 \sqrt{\|p_i\|^2-\rho^2}\,\mathsf Jp_i}{\|p_i\|^2}
 \quad(\|p_i\|\ge\rho).
\]
Thus one side of a least enclosing equilateral triangle either contains two
points, or is tangent to the disk and contains one point.
\end{proposition}

\begin{proof}
Apply \zcref{lem:support-cell-rotation}. A point--point tie has normal
perpendicular to $p_k-p_i$; solving $\langle p_i,n\rangle=\rho$ with
$\|n\|=1$ gives the displayed point--disk normals. Rotation through
$120$ degrees identifies the three supporting directions. A disk-only
cell contributes $3\rho$ and is the only additional alternative.
\end{proof}


\readerheading{How to use the finite list.}
For points, evaluate the three-support sum at the listed pair normals.
For a disk and points, also test the point--disk ties. Nonexposed ties
may be discarded, but their completeness must first be justified.
The disk-only alternative contributes $3\rho$ only at an orientation
where all three disk supports dominate every point; it is not an
unconditional candidate value. Equivalently, adjoining one arbitrary
reference normal to the complete tie list and evaluating the actual
support sum at every listed normal gives a finite test, including the
case with no ties. The final contact description in the proposition
applies to its non-disk-only alternative.

\readerheading{One disk and one point.}
A useful explicit specialization, with $r>0$ and $d=\|q\|$, is
\[
\Lambda(\mathbb B(O,r)\cup\{q\})=
\begin{cases}
2\sqrt3\,r,&d\le2r,\\
\sqrt3\,r+\sqrt{d^2-r^2},&d>2r.
\end{cases}
\]
Indeed, when $d\le2r$, orient the point projections as
$d/2,d/2,-d$, so all three disk supports suffice. When $d>2r$,
the support-cell argument puts a minimum at a point--disk tie.
If one projection equals $r$, the other two are
$(-r\pm\sqrt{3(d^2-r^2)})/2$; substitution in the support sum
gives the second formula. A circle gives the same enclosure problem
as its closed disk, since triangles are convex.
This example explains the transition but does not replace the
disk--finite-set proposition: several points can impose a joint
obstruction that none imposes alone.

\readerheading{The other structural step.}
The fixed radial-forcing lemma and common-pair domination below
turn local reach bounds into points missed by every open V triangle.
The uniform common-pair corollary then supplies a symmetric inner
hexagon and its inscribed disk. These are witness-construction results,
not consequences of the finite caliper criterion. BC and D subsequently
use fixed point sets with no disk; F uses nine fixed points and a
disk-plus-three-point comparison set.

\subsection{Local capacities and radial witnesses}

The following capacities answer a local question: after boundary anchors
have been prescribed, how far toward the center can a unit triangle reach?
They are maxima over a family of triangles, not the actual reaches of the
original triangle. A larger prescribed anchor shrinks the feasible family.
The neighboring capacities are needed because a point on one radial segment
may be covered by the V triangle at either adjacent vertex.

Normalize at \(V_0\) and define the four-anchor hull
\[
 K(a,b,c):=
 \conv\{V_0,X_5(1-a),X_0(b),(1-c)V_0\},
\]
\[
 \mathcal A:=
 \{(a,b,c)\in[0,1]^3:\Lambda(K(a,b,c))\le1\}.
\]
For
\[
 0\le a,b\le1,
 \qquad
 a^2+ab+b^2\le1,
\]
put
\[
 c_{\max}(a,b):=
 \max\{c\in[0,1]:(a,b,c)\in\mathcal A\}.
\]
For \(0\le a,c\le1\), also put
\[
 M_c(a):=\max\{b\in[0,1]:(a,b,c)\in\mathcal A\}.
\]
In particular, the \emph{diameter envelope} is
\begin{equation}
 M_0(x)=\frac{-x+\sqrt{4-3x^2}}2\qquad(0\le x\le1).
 \label{eq:diameter-envelope-definition}
\end{equation}
Here $M_0$ without an argument denotes the radial midpoint, whereas
$M_0(x)$ denotes this scalar function. For later rescuer estimates define
\begin{equation}
 M_c^{\rm sup}:=\frac{c+\sqrt{c^2-8c+4}}2\qquad(0\le c\le1/2).
 \label{eq:supercritical-envelope-definition}
\end{equation}
The value at $c=1/2$ is a continuous endpoint convention; no
supercritical triangle has that radial demand. The bound associated with
this envelope is proved in \zcref{lem:strict-supercritical-envelope}.
The exact admissible set and radial envelope are proved in
\zcref{prop:exact-admissible-set}.

The neighboring-arm capacities are the variational quantities
\[
 C_+(a,b):=
 \max\left\{c\in[0,1]:
 \begin{array}{l}
 \exists T\ [T\text{ is a closed unit equilateral triangle},\\
 \{V_0,X_5(1-a),X_0(b),(1-c)V_1\}\subseteq T]
 \end{array}\right\},
\]
\[
 C_-(a,b):=C_+(b,a),
 \qquad
 0\le a,b\le1,
 \qquad
 a+b\le1.
\]
Only a bound, not the full piecewise neighboring capacity, is needed below.
The direct four-caliper proof is \zcref{lem:direct-neighbor-diagonal}.

\begin{lemma}[Common-pair domination]
\label{lem:new-common-pair-domination}
If
\[
 p,q\ge0,
 \qquad
 p+q\le1,
\]
then
\[
 C_+(p,q),C_-(p,q)
 \le1-\min\{p,q\}\le c_{\max}(p,q).
\]
\end{lemma}

\begin{proof}
This is the direct diagonal-capacity comparison in
\zcref{lem:appendix-common-pair-domination}.
\end{proof}

\begin{definition}[Fixed total radial endpoint]
\label{def:fixed-total-radial-endpoint}
Parametrize $r_i$ by $Z_i(c)=(1-c)V_i$, measured from $V_i$ toward $O$.
Let $u_{j\to i}$ be the largest parameter in the actual positive trace of
$T_j$ on $r_i$ for $j=i-1,i+1$, with value zero when no such positive trace
exists. Put
\[
 \gamma_i:=\max\{C_i,u_{i-1\to i},u_{i+1\to i}\},\qquad
 d_i:=1-\gamma_i,\qquad \widehat P_i:=d_iV_i.
\]
\end{definition}

Thus $\widehat P_i$ is the innermost endpoint reached by any local V
triangle on $r_i$, rather than just the endpoint of the own V triangle.
This distinction is essential when a nonsupercritical V triangle supports
an adjacent radial segment.

\begin{lemma}[Fixed radial forcing and demand recovery]
\label{lem:fixed-total-radial-forcing}
One has $0<d_i<1$ and $\widehat P_i\notin\bigcup_jU_j$. Under skeleton
coverage, $\widehat P_i\in U_C$. More generally, for every
$\gamma_i\le\Gamma\le1$, the point $(1-\Gamma)V_i$ is missed by every
open V triangle and is center-forced under skeleton coverage.
If an arbitrary open unit candidate $U'$
contains $O$ and $\widehat P_i$, and its exit on $r_i$ is $d_i'$ measured
from $O$, then $\gamma_i>1-d_i'$. If both neighboring endpoints are at most
$1-d_i'$, it follows that
\[
 C_i>1-d_i'.
\]
It is sufficient that each neighboring V triangle be nonsupercritical with
both actual boundary reaches greater than $d_i'$.
\end{lemma}

\begin{proof}
A local open triangle containing the total endpoint would extend its trace
farther toward $O$. Nonlocal triangles are excluded by diameter. The last
assertion follows from common-pair domination, which bounds a neighboring
endpoint by $1-\min\{A_j,B_j\}$. The full endpoint and strictness argument
is in \zcref{lem:appendix-fixed-total-endpoint}.
\end{proof}

Let selected lower bounds satisfy
\[
 0\le a_j\le A_j,
 \qquad
 0\le b_j\le B_j
 \qquad(0\le j\le5),
\]
\[
 a_j^2+a_jb_j+b_j^2\le1
 \qquad(0\le j\le5).
\]
A neighboring term is used only when the actual V triangle has the
corresponding positive radial trace:
\[
 \Haus(T_{i-1}\cap r_i)>0
 \Longrightarrow a_{i-1}+b_{i-1}\le1,
\]
\[
 \Haus(T_{i+1}\cap r_i)>0
 \Longrightarrow a_{i+1}+b_{i+1}\le1.
\]
Define
\[
 \Gamma_i:=\max\Bigl(
 \{c_{\max}(a_i,b_i)\}
 \cup\{C_+(a_{i-1},b_{i-1}):\Haus(T_{i-1}\cap r_i)>0\}
\]
\[
 {}\hspace{35mm}\cup
 \{C_-(a_{i+1},b_{i+1}):\Haus(T_{i+1}\cap r_i)>0\}
 \Bigr),
 \qquad
 D_i:=(1-\Gamma_i)V_i.
\]
The points \(D_i\) are the radial witnesses used below.

\readerheading{Actual endpoint versus capacity bound.}
The number $\gamma_i$ records the innermost actual V-triangle trace on
$r_i$; $\widehat P_i=(1-\gamma_i)V_i$ is its endpoint. The number $\Gamma_i$
is an upper bound computed from selected demands, so $\Gamma_i\ge\gamma_i$.
Consequently $D_i=(1-\Gamma_i)V_i$ lies at least as far toward the center.
Both points are missed by the open V triangles, but for different reasons:
$\widehat P_i$ uses actual maximality, while $D_i$ uses a bound valid for every
allowed local triangle. Using only the own reach $C_i$ would overlook a
neighboring triangle's positive radial trace.


\begin{lemma}[Type-aware radial forcing]
\label{lem:new-type-aware-radial-forcing}
For every \(i\),
\[
 D_i\in r_i\subseteq S,
 \qquad
 D_i\notin\bigcup_{j=0}^5U_j.
\]
Consequently,
\[
 S\subseteq U_C\cup\bigcup_{j=0}^5U_j
 \Longrightarrow D_i\in U_C.
\]
\end{lemma}

\begin{proof}
The variational capacities give $\gamma_i\le\Gamma_i\le1$, including
only actual positive adjacent traces. Apply
\zcref{lem:fixed-total-radial-forcing}. The point $O$ when $\Gamma_i=1$
is handled by that lemma as well; no positive-distance claim is made there.
\end{proof}



\begin{corollary}[Uniform common-pair radial forcing]
\label{cor:new-uniform-common-pair-forcing}
Suppose $p,q\ge0$, $p+q\le1$, and
\[
 p\le a_i\le A_i,\qquad q\le b_i\le B_i\quad(0\le i\le5),
 \qquad 0<c_*:=c_{\max}(p,q)<1.
\]
Then $(1-c_*)V_i\notin\bigcup_{j=0}^5U_j$ for every $i$, without a
restriction on the normalized V types. Under a hypothetical cover of $H$,
\[
 (1-c_*)V_i\in U_C\quad(0\le i\le5),
 \qquad \mathbb B(O,h(1-c_*))\subset U_C.
\]
\end{corollary}

\begin{proof}
Reselect each boundary pair as $(p,q)$, using convexity.
Lemma~\ref{lem:new-type-aware-radial-forcing} applies to these common bounds,
including only the actual positive neighboring traces.
By \zcref{lem:new-common-pair-domination}, every included $C_+$ or $C_-$
term is at most $c_*$, while the own-ray term equals $c_*$. Hence each
maximum is exactly $c_*$ and the forced points are the stated symmetric
ones. Their convex hull contains the displayed disk. In particular, a
Vd1, Vd2, or T3-like crossing does not change any witness coordinate.
\end{proof}


% Active nonzero-gap body. Included by 06_finite_enclosure_full.tex.
% Witness sets are fixed before any enclosing candidate is introduced.
\subsection{Fixed witnesses and candidate enclosures}
\label{subsec:trace-exact-ab-envelopes}

Each witness set is fixed from the original V triangles before an
arbitrary enclosing candidate is chosen. The finite calipers vary their support orientation, but the witness points do not. Gap endpoints are supplied
by \zcref{lem:gap-exhaustion}; radial witnesses use the total endpoints,
including all permitted neighboring traces.







Replacement is a separate skeleton-preserving operation, not a fourth
witness construction. BC handles the center-aligned case. For an off-center
supercritical role, its midpoint supplier leads to D, replacement, or the
Vd2 perimeter obstruction.

\subsection{A selected gap forces five points (BC)}

Normalize the center midpoint to $M_0$ and choose any actual gap, reflected
to $J_0$. The other incident edge $e_{5,0}$ may or may not contain a gap.
Only the four middle edges $e_{1,2},\ldots,e_{4,5}$ must be gap-free.
The following transfer supplies the tail bound needed by the path argument.

\begin{lemma}[Shared-gap-anchor boundary transfer]
\label{lem:shared-gap-anchor-transfer}
If the original open roles cover $\partial H$ and $J_0$ is an actual gap,
then, independently of the gap status of $e_{5,0}$,
\[
 B_5>1-M_0(B_0)>B_0/2.
\]
Here $M_0(z)$ is the scalar diameter envelope.
\end{lemma}
\begin{proof}
Put $G=X_0(B_0)$ and $L=X_5(B_5)$. Actual gap forcing gives
$G\in U_C$, while maximality gives $G\in T_0$ and $L\notin U_5$.
Boundary locality and the original perimeter cover give
$L\in U_0\cup U_C$. Thus one unit triangle contains $L$ openly and
$G$ in its closure, so $\|G-L\|<1$. Indeed, at equality one could move
$L$ slightly away from $G$ inside that open triangle and exceed its diameter.
The two incident-edge directions at $V_0$ meet at $120$ degrees, hence
\[
 B_0^2+B_0(1-B_5)+(1-B_5)^2<1.
\]
Apply \zcref{lem:signed-diameter-transfer}; its last inequality is strict
because $B_0>0$. The proof includes singleton gaps and does not require
$L$ to be a witness in the final finite set.
\end{proof}

In the pure enclosure statement below the tail bound $B_5\ge B_0/2$
is explicit. In the covering application, the preceding lemma establishes
it for the original roles before an arbitrary enclosing candidate is chosen.
We do not assume that the candidate also covers the discarded gap.

Define the two nonuniform comparison pairs and radii by
\[
 x=B_0,\qquad y=1-A_1,\qquad z=B_3.
\]
\[
 r=1-c_{\max}(1-y,z),\qquad t=1-c_{\max}(1-z,x/2).
\]
Put $\widehat t=\min\{t,A_3\}$ and fix
\[
 K_{BC}:=\{M_0,X_0(x),X_0(y),rV_2,\widehat tV_4\}.
\]
There are at most five points, including singleton gaps. The midpoint is a
structural anchor; the other points are missed by every open V role.
The cap by $A_3$
controls the remaining neighboring supplier on $r_4$; it is not silently dropped.

\readerheading{The five-point selected-gap construction.}
Set $x=B_0$, $y=1-A_1$, and $z=B_3$. The two boundary pairs
$(1-y,z)$ and $(1-z,x/2)$ are different; they preserve the shared boundary
parameter $z$ instead of replacing all path data by a uniform minimum.
Their capacity deficits give $r$ and $t$. The five witnesses are
$M_0,X_0(x),X_0(y),rV_2,\min(t,A_3)V_4$.
There is no point on $r_3$, although $B_3$ remains part of the construction.
The cap by $A_3$ handles a neighboring supplier. When it is active,
a pair of witnesses already has distance greater than one; otherwise the capacity inequality verifies the single decisive caliper
of the five-point threshold criterion. Singleton gaps merely identify
$X_0(x)$ and $X_0(y)$.


\begin{theorem}[Unified five-point selected-gap enclosure]
\label{prop:new-nplus-one-all-vd0}
\label{lem:appendix-fixed-transverse}
Suppose $J_0$ is an actual gap, $A_i+B_i\le1$ for $1\le i\le5$, and
\[
 B_i+A_{i+1}>1\quad(1\le i\le4),\qquad B_5\ge B_0/2.
\]
Then $\Lambda(K_{BC})>1$. Under skeleton coverage with distinguished
C midpoint $M_0$, the fixed set $K_{BC}$ lies in $U_C$.
There is no condition on the gap status of $e_{5,0}$ or supercriticality
of $T_0$, and no type restriction on the nonsupercritical path.
\end{theorem}
\begin{proof}
The boundary-deficit identities give increasing $A_1,\ldots,A_5$ and
decreasing $B_1,\ldots,B_5$. Thus $0<x\le y<1$ and $x/2\le z\le y$.
Put $\rho_i=f(A_i,B_i)$. By \zcref{cor:clipped-radial-forcing},
the safe radii are $\min\{\rho_i,A_{i-1},B_{i+1}\}$ for $i=2,3,4$.
The bounds $(A_2,B_2)\ge(1-y,z)$ imply
$r\le\rho_2$, $r\le A_1$, and $r\le B_3$. Hence $rV_2$ is safe.
Similarly $(A_4,B_4)\ge(1-z,x/2)$ gives $t\le\rho_4$, and
$t\le x/2\le B_5$; clipping by $A_3$ makes $\widehat tV_4$ safe.
The gap endpoints follow from \zcref{lem:gap-exhaustion}, and the midpoint
belongs to $U_C$ by hypothesis. All witnesses are fixed before enclosure is tested.

If $A_3\ge t$, apply the capacity instance of \zcref{lem:bc-five-point}. Otherwise $A_3<t$ and
$A_3\ge A_1=1-y$, so
\[
 \|X_0(y)-A_3V_4\|^2\ge
 \|X_0(y)-(1-y)V_4\|^2=1+(1-y)(2-y)>1.
\]
The diameter obstruction gives $\Lambda(K_{BC})>1$ in this case also.
No signed center coordinates or CE1 return are used in this enclosure argument.
\end{proof}

\begin{theorem}[Center-aligned path obstruction]
\label{thm:center-aligned-path}
Suppose a skeleton cover has at least one actual gap and its C triangle's
unique midpoint is $M_k$. If every $T_i$ with $i\ne k$ is nonsupercritical,
the configuration is impossible.
\end{theorem}
\begin{proof}
Normalize $k=0$. The signed normal form confines possible gap edges to
$e_{5,0},e_{0,1}$. Select either actual gap and reflect it to $J_0$.
The four middle edges are gap-free, and
\zcref{lem:shared-gap-anchor-transfer} gives $B_5>B_0/2$.
Apply \zcref{prop:new-nplus-one-all-vd0}. The midpoint belongs openly
to $U_C$, and all other witnesses are center-forced. Compact-open
containment contradicts $\Lambda(K_{BC})\ge1$ in either gap rank.
\end{proof}

\begin{theorem}[Every skeleton cover has a supercritical V triangle]
\label{prop:new-nplus-zero-gap-closures}
\label{thm:fixed-nzero}
A skeleton cover by the seven original open roles satisfies $N_+\ge1$.
\end{theorem}
\begin{proof}
If $N_+=0$ and there are no gaps, strict overlaps give
$6<\sum_i(A_i+B_i)\le6$. If a gap is present, all six roles are
nonsupercritical, so \zcref{thm:center-aligned-path} applies. This proof uses only BC; replacement will invoke N0, not conversely.
\end{proof}

\begin{lemma}[A unique midpoint supplier]
\label{lem:midpoint-supplier}
In a nonzero-gap skeleton cover, the preceding N0 theorem and the skeleton
budget leave only $N_+=1$, $N_{\rm sp}\le1$. Write $M_k$ for the C
midpoint and $\sigma$ for the supercritical index. If $\sigma\ne k$,
there is a unique positive-support role $T_\tau$, where
\[
 \tau\in\{\sigma-1,\sigma+1\},\qquad M_\sigma\in U_\tau.
\]
If $T_\tau$ is T3-like, then $\tau=k$. If $\tau\ne k$, the edge shared
by $T_\tau,T_\sigma$ is center-free, in both CE1 and CE2.
\end{lemma}
\begin{proof}
N0 excludes $N_+=0$. Two supercritical roles would require a distinct
positive-support rescuer by \zcref{lem:positive-support-rescuer}, and
\zcref{thm:common-skeleton-count} would exclude them. The same count
bound gives $N_{\rm sp}\le1$.

If $\sigma\ne k$, neither the C triangle nor $T_\sigma$ contains
$M_\sigma$. Diameter excludes nonlocal roles; a Vd0 role has no positive
adjacent trace and cannot contain that point openly. Hence its unique
supplier is an adjacent positive-support role. A T3-like supplier also
misses its own midpoint by \zcref{lem:t3-nonsupercritical}. Unless $\tau=k$, neither the center nor the
other roles could cover this second midpoint: all other roles are Vd0.
Finally, when $k\notin\{\tau,\sigma\}$, the shared edge is not incident
to $V_k$, whereas every open center boundary trace is incident to $V_k$
by the signed normal form.
\end{proof}

\subsection{A supported midpoint forces four points (D)}

Let $Y(t)=(1-t)V_0+tV_5=X_5(1-t)$. The geometric lemma below is stated
independently of any V type or covering configuration.

\begin{theorem}[Four-point rescuer geometry]
\label{thm:fixed-four-point-rescuer}
\label{lem:appendix-fixed-four-point}
Suppose $a\ge0$, $\varepsilon>0$, $\beta\ge0$, and
\[
 s=a+\varepsilon,\qquad \beta\le\frac{\varepsilon}{s}.
\]
Then
\[
 \Lambda\{O,\varepsilon V_1,Y(a),Y(1-\beta)\}\ge1.
\]
Neither $a\le\varepsilon$ nor $s\le1$ is required by this geometric statement.
\end{theorem}
\begin{proof}
If $a=0$, the set contains $O,V_0$, at distance one. If $s\ge1$, then
\[
 \|Y(a)-\varepsilon V_1\|^2=1-s+s^2\ge1.
\]
These are diameter obstructions. Otherwise put $v=a/s$, so $0<s,v<1$.
The ratio gives $1-\beta\ge v$, while $a=sv\le v$. Hence $Y(v)$ lies
on $[Y(a),Y(1-\beta)]$. By convexity it suffices to enclose the quadrilateral,
in cyclic order,
\[
 O,\qquad Q=Y(v),\qquad A=Y(sv),\qquad P=s(1-v)V_1.
\]
With $h=\sqrt3/2$, its outward edge normals may be chosen as
\[
 (-hv,v/2-1),\quad v(1-s)(h,-1/2),\quad
 (hs,1-s/2),\quad s(1-v)(-h,1/2).
\]
The maximizing triples in each direction and its two $120$-degree rotations
are $(O,A,P)$, $(A,P,O)$, $(A,O,Q)$, and $(O,Q,A)$, respectively.
The resulting caliper sides are
\[
 \frac1{\sqrt{1-v+v^2}},\qquad 1+s(1-v),\qquad
 \frac1{\sqrt{1-s+s^2}},\qquad 1+v(1-s).
\]
All exceed one for $0<s,v<1$. The finite-caliper theorem
\zcref{thm:cert-caliper} completes the nondegenerate case.
\end{proof}

\readerheading{Why the ratio conditions matter.}
The identity in the preceding proof combines the supported radial point
with one boundary point to produce a point on $r_0$. Since $a\le\varepsilon$,
its distance from $O$ is $\varepsilon/s\ge1/2$. A candidate also
contains $O$, so it must contain $M_0$. The other boundary point then demands
more of the candidate's left boundary trace than its signed-center geometry
allows. This is why the application needs a ratio estimate: it converts the
local supplier geometry into this same four-point contradiction.


\begin{lemma}[Scalar rescuer ratio test]
\label{lem:rescuer-ratio-test}
For $0\le x,c\le1/2$ and $M=(c+\sqrt{c^2-8c+4})/2$,
\[
 x\le1-M\quad\Longleftrightarrow\quad x^2+(c-2)x+c\ge0.
\]
\end{lemma}
\begin{proof}
The smaller root is $1-M$, while the other root is at least one.
The quadratic is nonnegative on $[0,1/2]$ exactly before its smaller root.
\end{proof}

For the covering application, let $T_0$ supply $M_1$, with supported
interval $[c,u]$ on $r_1$, measured from $V_1$ toward $O$, and put
$\varepsilon=1-u$ and $P_T=\varepsilon V_1$. All interval endpoints here
belong to the original triangle, not to its translated majorant.
The local T3-like and Vd1 calculations establish
\[
 A_0+\varepsilon\le1,\qquad
 \frac{A_0}{A_0+\varepsilon}\le1-M_c^{\rm sup}.
\]
Together with $C_1\ge c$, the boundary path, and the forced endpoint,
\zcref{lem:new-rescuer-tail-budget} gives the single fixed witness
\[
 K_D=\{O,\varepsilon V_1,Y(A_0),Y(1-B_5)\}\subset U_C,
 \qquad \Lambda(K_D)\ge1.
\]
The two local verifications are
\zcref{lem:appendix-t3-supported-tail,lem:appendix-vd1-supported-tail};
the enclosure step does not distinguish one from two gaps.



\begin{figure}[!htbp]
\centering
\begin{minipage}[t]{.47\linewidth}
\centering
\resizebox{.96\linewidth}{!}{%
\begin{tikzpicture}[x=6cm,y=1cm,>=Latex]
  \node[panellabel,anchor=west] at (-.02,2.30) {(a) the supported interval on $r_1$};
  \draw[flow] (0,1.15)--(1,1.15);
  \fill (0,1.15) circle (.017) node[left=2pt,figlabel] {$V_1$};
  \fill (1,1.15) circle (.017) node[right=2pt,figlabel] {$O$};
  \coordinate (C) at (.24,1.15);
  \coordinate (M) at (.50,1.15);
  \coordinate (U) at (.73,1.15);
  \draw[actualreach] (C)--(U);
  \draw[VertexOrange,fill=white,line width=.7pt] (U) circle (.025);
  \fill[DemandRed] (U) circle (.011);
  \fill[black] (C) circle (.019) node[above=3pt,figlabel] {$c$};
  \fill[CenterBlue] (M) rectangle +(0.026,0.026);
  \node[above=3pt,figlabel] at (M) {$M_1$};
  \node[above=3pt,figlabel] at (U) {$u$};
  \draw[dimline] ($(U)+(0,-.22)$)--($(1,1.15)+(0,-.22)$)
    node[midway,below=1pt,figlabel] {$1-u$};
  \node[figlabel,align=center] at (.52,.12)
    {$P_T=(1-u)V_1$ is the O-side endpoint\\
     coordinate on the top axis is measured from $V_1$ toward $O$};
\end{tikzpicture}%
}
\end{minipage}\hfill
\begin{minipage}[t]{.47\linewidth}
\centering
\resizebox{.96\linewidth}{!}{%
\begin{tikzpicture}[scale=2.15,>=Latex]
  \node[panellabel,anchor=west] at (-1.08,1.22) {(b) why $P_T$ is center-forced};
  \HexCoords
  \DrawHexagon
  \DrawRays
  \coordinate (PT) at ($(O)!.27!(V1)$);
  \coordinate (M1) at ($(O)!.5!(V1)$);
  \draw[vertexrole]
    (V0)--($(V0)!.76!(V1)$)--($(O)!.20!(V1)$)--cycle;
  \draw[DemandRed,dashed,line width=.9pt]
    (V1)--($(V1)!.48!(V0)$)--($(O)!.48!(V1)$)--cycle;
  \draw[actualreach] ($(V1)!.22!(O)$)--($(V1)!.73!(O)$);
  \fill[DemandRed] (PT) circle (.025) node[left=2pt,figlabel] {$P_T$};
  \fill[CenterBlue] (M1) rectangle +(0.025,0.025);
  \node[figlabel,right] at (M1) {$M_1$};
  \draw[flow] (PT)--($(PT)!1.55!(O)$);
  \node[figlabel,align=left,anchor=west] at (-.96,-.72)
    {$T_0$: open endpoint at $P_T$\\
     $T_1$: misses $M_1$, hence the O-side segment\\
     other V roles: excluded by type or diameter};
  \node[figlabel,atlasbox,align=center] at (.48,-1.05)
    {$P_T\in U_C$};
\end{tikzpicture}%
}
\end{minipage}
\caption{Construction of the T3-like rescuer endpoint.  The interval
$[c,u]$ is parametrized from $V_1$ toward $O$, so its O-side endpoint is
$P_T=(1-u)V_1$.  The T3-like open role omits its endpoint, the adjacent
supercritical role contains $V_1$ but misses $M_1$ and therefore misses the
O-side segment by convexity, and all remaining V roles are excluded.  Thus
$P_T$ is forced into the C triangle.}
\label{fig:t3-rescuer-endpoint}
\end{figure}

\FloatBarrier

\readerheading{A change of configuration, not a new candidate.}
The away-supplier case uses a different operation from finite enclosure.
Under the stated two-chart hypotheses, two V triangles are replaced while
the distinguished roles and coverage of the skeleton are retained. The
result in this unique-supercritical case has no supercritical V triangle, contradicting
\zcref{thm:fixed-nzero}. This explains why that theorem was proved for
\emph{skeleton} covers before replacement was introduced.

Neither coverage of all of $H$ nor the original number of gaps is claimed
to survive. The contradiction uses only the preserved skeleton cover.
Likewise, the earlier translated closed-trace majorant is not itself a
valid replacement open triangle: its distinguished vertex lies on its
boundary. These three objects---a majorant, an enclosing candidate, and a
replacement arrangement---have different uses.



\subsection{The zero-gap nine-point obstruction}
\label{sec:strategy4-reader}

The remaining zero-gap case has one supercritical V triangle. There is no
boundary gap to use as a witness, so the obstruction must be built inside
$H$. The construction and its verification have three stages.

\readerheading{Force the points.}
Strict boundary handoffs give common lower bounds for the boundary reaches.
These force six radial points into the C triangle. Three further points are
chosen on the exact frontier of the feasible union of the supercritical
triangle; separate diameter and handoff arguments exclude every other V
triangle. Exclusion from the supercritical triangle alone would not be enough.
A frontier point may belong to a closed feasible triangle while being missed
by its open interior.

\readerheading{Simplify the set to be enclosed.}
The convex hull of the six radial points contains a disk. Replacing those
six points by this disk gives a smaller comparison set. In the nontrivial
range, the two outer frontier points are then moved inward by a single
Newton step at the fixed line parameter $k=1/2$. Convexity keeps these comparison points inside any candidate
containing the original witnesses. These are exact inner comparisons,
not numerical approximations to a covering configuration.

\readerheading{Bound the enclosing triangle.}
The inner comparison reduces the enclosure problem to four supporting
contacts: two contacts along consecutive-point lines and two disk-tangent
contacts. The line estimates are analytic. The remaining paired tangent
inequalities are reduced to the exact polynomial certificate. At the end,
the comparison set requires side at least one, contradicting its containment
in the open C triangle. The calculations below justify each of these stages;
the certificate enters only at the last one.

\readerheading{Active tangent calculation.}
The two outer comparison points use the fixed-start formulas, not the older
junction-start formulas. The exact coefficient comparisons and radius envelope
are specified in \zcref{thm:paper-exact-mixed-certificate}.
The actual disk is unchanged; the smaller scalar radius is used only for
initially checking the tangent residuals.


Assume
\[
 N_{\rm gap}=0,
 \qquad
 N_+=1.
\]
Rotate the unique supercritical role to index \(4\).  The exact-one handoff
calculation \zcref{lem:ab-extreme-jump} gives
\begin{equation}
 \xi_3<\xi_2<\xi_1<\xi_0<\xi_5<\xi_4.
 \label{eq:reader-extreme-chain}
\end{equation}
Define
\begin{equation}
 a:=1-\xi_3,\qquad
 b:=\xi_4,\qquad
 p:=1-b,\qquad
 q:=1-a.
 \label{eq:reader-ab-parameters}
\end{equation}
Then
\begin{equation}
 \begin{gathered}
 0<a,b<1,\qquad a=a_4<A_4,\qquad b=b_4<B_4,\\
 a+b>1,\qquad a^2+ab+b^2<1,\\
 a_i\ge p,\qquad b_i\ge q\quad(0\le i\le5).
 \end{gathered}
 \label{eq:reader-ab-domain}
\end{equation}

Define the two boundary anchors
\[
 P_A:=V_4+a(V_3-V_4),
 \qquad
 P_B:=V_4+b(V_5-V_4),
\]
and the feasible union
\[
 \mathcal R_{AB}(a,b):=
 \bigcup\left\{T\cap H:
 \begin{array}{l}
 T\text{ is a closed unit equilateral triangle},\\
 V_4,P_A,P_B\in T
 \end{array}\right\}.
\]
Put
\[
 c_*:=c_{\max}(p,q),
 \qquad
 r_*:=h(1-c_*),
 \qquad
 P_i^{\rm rad}:=(1-c_*)V_i\quad(0\le i\le5).
\]
The bounds in \zcref{lem:symmetric-core-witness} give $0<c_*<1$.
Applying \zcref{cor:new-uniform-common-pair-forcing} to the common pair gives
\begin{equation}
 \mathbb B(O,r_*)
 \subseteq\conv\{P_0^{\rm rad},\ldots,P_5^{\rm rad}\}
 \subset U_C.
 \label{eq:reader-forced-disk}
\end{equation}

The two circles adjacent to the supercritical frontier are
\[
 \mathcal C_-:=\partial\mathbb B(X_2(b),1),
 \qquad
 \mathcal C_+:=\partial\mathbb B(X_5(1-a),1).
\]
Let $\ell_-,\ell_+$ be the two straight pieces of the interior frontier
of $\mathcal R_{AB}(a,b)$, in the order of
\zcref{thm:strict-ab-union}. Define
\[
 \{Q_0\}=\ell_-\cap\ell_+,\qquad
 \{Q_-\}=\ell_-\cap\mathcal C_-,\qquad
 \{Q_+\}=\ell_+\cap\mathcal C_+.
\]
The coordinate formulas and selected first roots are given in
Appendix~\ref{sec:zero-gap-nine-point-optimization},
\zcref{eq:line-junction,eq:asym-witnesses}. Lemma~\ref{lem:fixed-line-signs}
proves that these intersections are singletons on the indicated pieces.
The nine-point set is
\[
 K_F:=\{P_i^{\rm rad}:0\le i\le5\}\cup\{Q_-,Q_0,Q_+\}.
\]

Their exact construction and exclusions in
\zcref{thm:strict-ab-union,lem:fixed-line-signs,
lem:asymmetric-core-witness} use contained vertices and actual handoffs, not
Vd0 locality, and give
\begin{equation}
 Q_-,Q_0,Q_+\in U_C.
 \label{eq:reader-forced-asymmetric}
\end{equation}

For the enclosure calculation use the compact comparison set
\begin{equation}
 K_{\rm wit}(a,b):=
 \mathbb B(O,r_*)\cup\{Q_-,Q_0,Q_+\}
 \subset U_C.
 \label{eq:reader-witness-set}
\end{equation}

\begin{figure}[!htbp]
\centering
\resizebox{.76\linewidth}{!}{%
\begin{tikzpicture}[scale=3.0,>=Latex]
  \HexCoords
  \DrawHexagon
  \DrawRays

  \foreach \i in {0,...,5}{
    \coordinate (Pr\i) at ($(O)!.16!(V\i)$);
    \fill[DemandRed] (Pr\i) circle (.018);
  }
  \draw[DemandRed!65,line width=.65pt]
    (Pr0)--(Pr1)--(Pr2)--(Pr3)--(Pr4)--(Pr5)--cycle;
  \draw[CenterBlue,dashed,fill=CenterBlue!7,line width=.75pt]
    (O) circle ({.16*sqrt(3)/2});

  \coordinate (Qm) at (-.18,-.50);
  \coordinate (Qz) at (-.02,-.62);
  \coordinate (Qp) at (.17,-.52);
  \fill[GapPurple] (Qm)--++(.025,.043)--++(.025,-.043)--cycle;
  \fill[GapPurple] (Qz)--++(.025,.043)--++(.025,-.043)--cycle;
  \fill[GapPurple] (Qp)--++(.025,.043)--++(.025,-.043)--cycle;
  \node[figlabel,left=2pt] at (Qm) {$Q_-$};
  \node[figlabel,below=2pt] at (Qz) {$Q_0$};
  \node[figlabel,right=2pt] at (Qp) {$Q_+$};
  \node[figlabel,right] at (Pr0) {$P_0^{\rm rad}$};
  \node[figlabel,above right] at (Pr1) {$P_1^{\rm rad}$};
  \node[figlabel,below left] at (Pr4) {$P_4^{\rm rad}$};
  \fill (O) circle (.012) node[above left=1pt,figlabel] {$O$};

  \draw[centerrole]
    (-.49,-.73)--(.50,-.72)--(.02,.24)--cycle;
  \node[figlabel,text=CenterBlue] at (.43,-.63) {putative $U_C$};

  \node[figlabel,atlasbox,align=left,anchor=west] at (1.20,.55)
    {six symmetric points: $P_i^{\rm rad}=(1-c_*)V_i$\\
     three asymmetric points: $Q_-,Q_0,Q_+$\\
     $\mathbb B(0,r_*)\subset\operatorname{conv}\{P_i^{\rm rad}\}$};
  \draw[flow] (1.15,.30)--(.34,.05);

  \node[figlabel,align=center] at (0,-1.22)
    {$K_{\rm wit}=\mathbb B(0,r_*)\cup\{Q_-,Q_0,Q_+\}
      \subset U_C$ under a hypothetical cover};
\end{tikzpicture}%
}
\caption{The zero-gap nine-point witness.  The six symmetric radial points
force a centered disk into the convex C triangle, while the three asymmetric
frontier points supply the missing directional support.  Marker shapes
separate the two constructions: circles are radial witnesses and triangles
are strict-supercritical $AB$-frontier witnesses.  The displayed placement is
illustrative; the point definitions are the exact ones in the preceding
lemmas.}
\label{fig:zero-gap-nine-point-witness}
\end{figure}


\readerheading{The witness set and the comparison set are different.}
The original finite witness is $K_F$, with at most nine points. The set
$K_{\rm wit}$ contains a disk in place of the six radial points and satisfies
\[
 K_{\rm wit}\subseteq\conv(K_F),\qquad
 \Lambda(K_F)=\Lambda(\conv(K_F))\ge\Lambda(K_{\rm wit}).
\]
Thus proving the lower bound for the disk comparison proves it for the
nine-point witness. The disk need not be missed by every V triangle:
its containment in $U_C$ follows from convexity and the six forced radial
points. The later Newton points play the same comparison role; they do not
replace the definitions of the frontier points $Q_-,Q_0,Q_+$.

\readerheading{Active tangent calculation.}
The two outer comparison points use the fixed-start formulas, not the older
junction-start formulas. The exact coefficient comparisons and radius envelope
are specified in \zcref{thm:paper-exact-mixed-certificate}.
The actual disk is unchanged; the smaller scalar radius is used only for
initially checking the tangent residuals.


\readerheading{Why four contacts suffice in the hard range.}
The additional geometric step is not another general enclosure
formula. The Newton inner reduction moves the two outer frontier
witnesses inward, preserving the candidate's containment by convexity.
Their convex hull with the disk has an exposed three-point chain;
its two successive point--point lines and two outer disk tangencies
supply the four contacts. The support-cell lemma then makes that
contact list exhaustive. The inner inclusion, contact geometry, and
tangent inequalities must still be proved; their calculations are in
\zcref{lem:technical-newton-reduction,prop:technical-four-contact-geometry,
lem:appendix-witness-enclosure-calculation}. Four contacts are not four
witness points, and the exact tangent certificate is not replaced by
the finite-candidate theorem.


\begin{theorem}[Witness enclosure]
\label{thm:reader-witness-enclosure}
For $0<a,b<1$, $a+b>1$, and $a^2+ab+b^2<1$, the comparison set
$K_{\rm wit}(a,b)$ satisfies
\[
 \Lambda(K_{\rm wit}(a,b))\ge1.
\]
\end{theorem}

\begin{proof}
The disk closes the range \(c_*\le2/3\).  For \(c_*>2/3\), the Newton inner
points reduce the enclosure to four supporting contacts: the lines through
two consecutive points and the two exposed disk tangencies. The line bounds
are analytic; the paired tangent residuals use the exact certificate and
simultaneous radius transfer.  The complete calculation is
\zcref{lem:appendix-witness-enclosure-calculation}.
\end{proof}

\begin{theorem}[Zero-gap obstruction]
\label{thm:reader-zero-gap-obstruction}
There is no hypothetical cover with
\[
 N_{\rm gap}=0,
 \qquad
 N_+=1.
\]
\end{theorem}

\begin{proof}
The construction gives the compact containment
\(K_{\rm wit}(a,b)\subset U_C\), while
\zcref{thm:reader-witness-enclosure} gives
\(\Lambda(K_{\rm wit}(a,b))\ge1\).  This contradicts
\zcref{lem:compact-cover-margin}.
\end{proof}

\subsection{Summary of the fixed witnesses}

\begingroup
\small
\renewcommand{\arraystretch}{1.18}
\setlength{\tabcolsep}{3pt}
\begin{longtable}{@{}p{.29\textwidth}p{.44\textwidth}p{.18\textwidth}@{}}
\caption{Fixed finite-enclosure recipes. Ordinary nonsupercritical V triangles
are not subdivided by type.}\label{tab:finite-enclosure-subcases}\\
\toprule
Conditions & Fixed witness set & Disk and count\\
\midrule
\endfirsthead
\toprule Conditions & Fixed witness set & Disk and count\\\midrule
\endhead
BC: at least one incident gap; $T_1,\ldots,T_5$ nonsupercritical &
$\{M_0,X_0(B_0),X_0(1-A_1),\widehat P_2,\widehat P_3,\widehat P_4\}$ & No disk; at most 6\\
D: center-based T3-like or Vd1 rescuer; one or two gaps &
$\{O,\varepsilon V_1,Y(A_0),Y(1-B_5)\}$ & No disk; at most 4\\
F: zero gaps, $N_+=1$ & The nine-point set $K_F$ & Disk; at most 9\\
\bottomrule
\end{longtable}
\endgroup
